\subsection{Capacitive Voltage Divider}

This block is constructed from two capacitors in series and is used
to attenuate a voltage signal.

\subsubsection{Schematic}
\begin{circuitfig}
  % Paths, nodes and wires:
  \node[ground, rotate=-90] at (5.5, 5.25){};
  \node[ground, rotate=-90] at (5.5, 4.75){};
  \draw (5.5, 6.75) to[sinusoidal voltage source, l={$V_{cc}$}] (5.5, 5.25);
  \draw (5.5, 4.75) to[sinusoidal voltage source, l={$V_{ee}$}] (5.5, 3.25);
  \draw (7, 5) -- (7.75, 5);
  \draw (7, 7) -| (5.5, 7) -| (5.5, 6.75);
  \node[shape=rectangle, minimum width=0.715cm, minimum height=0.465cm](N1) at (8.125, 5){} node[anchor=center] at (N1.text){$V_{out}$};
  \draw (7, 5.25) -| (7, 5);
  \draw (7, 4.75) -| (7, 5);
  \draw (5.5, 3.25) -| (5.5, 3) -- (7, 3);
  \draw (7, 7) to[capacitor, l={$C_1$}] (7, 5.25);
  \draw (7, 4.75) to[capacitor, l={$C_1$}] (7, 3);
\end{circuitfig}

\subsubsection{Transfer Function}
Applying KCL at $V_{out}$ we get:
\[
  (V_{out} - V_{cc}){sC_1} + (V_{out} - V_{ee}){sC_2} = 0
\]
\paragraph{Superposition:}
\begin{enumerate}
  \item Due to \(V_{cc}\) alone (\(V_{ee} = 0V\)):
    \begin{align*}
      (V_{out} - V_{cc}){sC_1} + (V_{out} - 0){sC_2} &= 0 \\
      (V_{out} - V_{cc}){sC_1} + V_{out}{sC_2} &= 0 \\
      V_{out}sC_1 - V_{cc}sC_1 + V_{out}sC_2 &= 0 \\
      V_{out}(sC_1 + sC_2) - V_{cc}sC_1 &= 0 \\
      V_{out}(sC_1 + sC_2) &= V_{cc}sC_1 \\
      \frac{V_{out}}{V_{cc}} &= \frac{sC_1}{sC_1 + sC_2} \\
      \frac{V_{out}}{V_{cc}} &= \frac{sC_1}{s(C_1 + C_2)} \\
      \frac{V_{out}}{V_{cc}} &= \frac{C_1}{C_1 + C_2} \\
      H_1(s) = \frac{N_1(s)}{D_1(s)} = \frac{V_{out_1}}{V_{cc}} = \frac{C_1}{C_1 + C_2}
    \end{align*}
  \item Due to \(V_{ee}\) alone (\(V_{cc} = 0V\)):
    \begin{align*}
      (V_{out} - 0){sC_1} + (V_{out} - V_{ee}){sC_2} &= 0 \\
      V_{out}sC_1 + (V_{out} - V_{ee}){sC_2} &= 0 \\
      V_{out}sC_1 + V_{out}sC_2 - V_{ee}sC_2 &= 0 \\
      V_{out}(sC_1 + sC_2) - V_{ee}sC_2 &= 0 \\
      V_{out}(sC_1 + sC_2) &= V_{ee}sC_2 \\
      \frac{V_{out}}{V_{ee}} &= \frac{sC_2}{sC_1 + sC_2} \\
      \frac{V_{out}}{V_{ee}} &= \frac{sC_2}{s(C_1 + C_2)} \\
      \frac{V_{out}}{V_{ee}} &= \frac{C_2}{C_1 + C_2} \\
      H_2(s) = \frac{N_2(s)}{D_2(s)} = \frac{V_{out_2}}{V_{ee}} = \frac{C_2}{C_1 + C_2}
    \end{align*}
\end{enumerate}

\paragraph{Total Output:}
\begin{align*}
  V_{out} &= V_{out_1} + V_{out_2} \\
  V_{out} &= V_{cc} \frac{C_1}{C_1 + C_2} + V_{ee} \frac{C_2}{C_1 + C_2}
\end{align*}

\subsubsection{Analysis}
\begin{enumerate}
  \item \textbf{Capacitor Ratio (\(\frac{C_1}{C_2}\)):}
    From KCL at \(V_{out}\):
    \begin{align*}
      (V_{out} - V_{cc})sC_1 + (V_{out} - V_{ee})sC_2 &= 0 \\
      (V_{out} - V_{cc})sC_1 &= -(V_{out} - V_{ee})sC_2 \\
      C_1(V_{out} - V_{cc}) &= -C_2(V_{out} - V_{ee}) \\
      \frac{C_1}{C_2} &= -\frac{V_{out} - V_{ee}}{V_{out} - V_{cc}}
    \end{align*}
  \item \textbf{Impedance's:}
    \paragraph{At Input:}
      \begin{align*}
        Z_{in} &= \frac{1}{sC_1} + \frac{1}{sC_2} \\
        Z_{in} &= \frac{1}{s}\left(\frac{1}{C_1} + \frac{1}{C_2}\right) \\
        Z_{in} &= \frac{1}{s}\left(\frac{C_1 + C_2}{C_1 C_2}\right) \\
        Z_{in} = \frac{C_1 + C_2}{s C_1 C_2} = \frac{1}{sC_{eq}} \hspace{1cm} \text{(where \(C_{eq} = \frac{C_1 C_2}{C_1 + C_2}\))}
      \end{align*}
    \paragraph{At Output:}
      \begin{align*}
        Z_{out} &= \frac{1}{sC_1 + sC_2} \\
        Z_{out} &= \frac{1}{s(C_1 + C_2)} \\
        Z_{out} = \frac{1}{s(C_1 + C_2)} = \frac{1}{sC_{eq}} \hspace{1cm} \text{(where \(C_{eq} = C_1 + C_2\))}
      \end{align*}
  \item \textbf{Loading Effect:} \\
    Similar to resistive dividers.
\end{enumerate}

\subsubsection{Question}
Design a capacitive voltage divider to attenuate 100V to 5V.
Assume \(C_1 = 10nF\) and \(V_{ee} = 0V\). Calculate the input
and output impedance of the divider.
\paragraph{Solution:}
\begin{enumerate}
  \item Calculate the capacitor ratio:
    \begin{align*}
      \frac{C_1}{C_2} &= -\frac{V_{out} - V_{ee}}{V_{out} - V_{cc}} \\
      \frac{C_1}{C_2} &= -\frac{5V - 0V}{5V - 100V} \\
      \frac{C_1}{C_2} &= 0.05263
    \end{align*}
  \item Calculate \(C_2\):
    \begin{align*}
      C_2 &= \frac{C_1}{0.05263} \\
      C_2 &= \frac{10nF}{0.05263} \\
      C_2 &= 190nF
    \end{align*}
  \item Calculate the input impedance:
    \begin{align*}
      Z_{in} &= \frac{C_1 + C_2}{s C_1 C_2} \\
      Z_{in} &= \frac{10nF + 190nF}{s (10nF)(190nF)} \\
      Z_{in} &= \frac{200nF}{s (1900nF^2)} \\
      Z_{in} &= \frac{2}{s 19nF} \\
      Z_{in} &= \frac{2}{j 2 \pi f 19nF} \hspace{1cm} \text{(substituting \(s = j \omega = j 2 \pi f\))}
    \end{align*}
  \item Calculate the output impedance:
    \begin{align*}
      Z_{out} &= \frac{1}{s(C_1 + C_2)} \\
      Z_{out} &= \frac{1}{s(10nF + 190nF)} \\
      Z_{out} &= \frac{1}{s(200nF)} \\
      Z_{out} &= \frac{1}{j 2 \pi f (200nF)} \hspace{1cm} \text{(substituting \(s = j \omega = j 2 \pi f\))}
    \end{align*}
\end{enumerate}

